\section{Symmetric dependent types}
\seclabel{dependent types}

To precisely show how our construction captures traditional \WLOG
arguments, we need to move to a more expressive setting. For that, we
use dependent types, which also allow us to impose richer invariants
on programming examples such as binary search trees.

We start by recalling the framework of \emph{families}
(\subsecref{families}), which gives a mathematical interpretation of
familiar dependent-type notions. Afterwards, we generalise our
constructs and results to a dependent setting (\subsecref{indexed
actions}) and show how via the Curry-Howard correspondence, they
capture proofs that use \WLOG arguments (\subsecref{propositional
reasoning}). Finally, we use the additional power of dependent types
to define a more refined notion of canoniser equivariance
(\subsecref{equivariant canonisers}).

\subsection{Families}
\subseclabel{families}

A \emph{family}~$\Ctx \types X$ is a pair consisting of a set $\Ctx$
and an assignment of a set $X \ctx$ to each element $\ctx \in \Ctx$.
We call $\Ctx$ the \emph{indexing} set, its elements $\ctx \in \Ctx$
\emph{indices}, and each set $X \ctx$ the \emph{fibre} over the index
$\ctx$. A \emph{term}~$\Ctx \types t : X$ in a family $\Ctx \types X$
is an assignment of an element $t \ctx \in X \ctx$ to each $\ctx \in
\Ctx$. A family $\Ctx \types X$ is a \emph{subfamily} of $\Ctx \types
Y$, which we write as $\Ctx \types X \subseteq Y$ if $X \ctx \subseteq
Y \ctx$ for all $\ctx \in \Ctx$.

\begin{example}
  \examplelabel{families example}
  To eliminate the necessity for the $\Maybe$ type in the tree lookup
  from \exampleref{tree example}, let us define the family $\Tree
  \types \NodeInName$ of \emph{node paths} that contains all paths
  ending at an internal node of the index tree. More precisely, the
  fibre $\NodeIn \btree$ is empty if $\btree$ is empty, and is
  otherwise the set of lists over $\TreeLabel = \set{\Lft,\Rgt}$ which are either
  $\Nil$, representing the root of the index tree, or $d \:: p$ that
  start with a step direction $d$ followed by a path $p$ in that
  subtree. This set can be given inductively as:
  \begin{gather*}
    \begin{array}[b]{c}
    \inferrule{
      ~
    }{
      \Nil \in \NodeIn \Big( \NodeV{\ell}x{r}\Big)
    }\quad
    \inferrule{
      p \in \NodeIn \ell
    }{
      (\Lft \:: p) \in \NodeIn \Big(\NodeV{\ell}x{r}\Big)
    }\quad
    \inferrule{
      p \in \NodeIn r
    }{
      (\Rgt \:: p) \in \NodeIn \Big(\NodeV{\ell}x{r}\Big)
    }
    \end{array}
    \tag*{\qedhere}
  \end{gather*}
\end{example}

When specifying fibres or terms as assignments of the form $\fun\mctx
X$ or $\fun\mctx t$, which explicitly use a generic index $\mctx
\in \Ctx$, we write them as $\mctx : \Ctx \types X$ and $\mctx : \Ctx
\types t : X$.
We will continue to distinguish these uses, denoting generic indices
for the purpose of defining a family by $\mctx : \Ctx$, and use $\mctx
\in \Ctx$ when we refer to elements of an indexing set $\Ctx$ for any other
purpose. Note that despite using a convenient type-theoretic notation,
we continue to keep our development purely mathematical, constructing
objects without any recourse to a syntactic type system.  Were we to
present such a syntactic type system and differentiate between
syntactic contexts $\Ctx$ and their interpretations as indexing sets
$\sem\Ctx$, then we would distinguish between $\ctx : \Ctx$ and $\ctx
\in \sem\Ctx$.


Continuing this approach, we allow fibres and terms to depend on a
context of variables. For the empty context, we take the indexing set
$\terminal \definedby \set\unitval$ and identify families $\terminal
\types X$ with their only fibre~$X \unitval$. Then, we write fibres
over $\terminal$ simply as $\types X$, and similarly for terms.

For context extension, we employ the \emph{Grothendieck construction}.
Recall that for a family $\Ctx \types X$, it is defined as the set
$\coprod_{\Ctx}X \definedby \coprod_{\ctx \in \Ctx} X \ctx$,
consisting of pairs $\pair\ctx x$ where $\ctx \in \Ctx$ and $x \in X
\ctx$, together with a projection function $\projection_1
\definedby \parent{\fun{\pair\ctx x} \ctx}: \coprod_{\Ctx}X \to \Ctx$.
We write $\Ctx, \mx : X \types Y$ and $\Ctx, \mx : X \types t : Y$ for
fibres and terms indexed by $\coprod_{\Ctx}X$ and use $\mx$ to refer
to the second component of the index pair from $\coprod_{\Ctx}X$ in $Y$
and $t$.

Given two families $\Ctx \types \mX$ and $\Ctx, \x : \mX \types \mY$,
we have the family of \emph{dependent pairs} $\Ctx \types
\DPair\x\mX\mY$ whose fibre at $\ctx \in \Ctx$ is the Grothendieck
construction $\coprod_{\mx \in \mX \ctx}\mY \pair\ctx\mx$, and the
family of \emph{dependent functions} $\Ctx \types \DFun\x\mX\mY$ whose
fibre at $\ctx$ is the product $\prod_{\mx \in \mX \ctx}\mY
\pair\ctx\mx$. As usual, we write $\Fun\mX\mY$ when the codomain
family $\mY$ does not depend on the argument of the function. To
distinguish between indices and arguments of dependent functions $\Ctx
\types f : \DFun\x\mX\mY$, we apply them as $f_{\ctx} \mx$
for $\ctx \in \Ctx$ and $\mx \in \mX \ctx$, or even as $f \mx$
when $\ctx$ can be inferred.

\begin{example}
  \examplelabel{node lookup}
  With dependent types, we can refine the lookup function $\Tree \to
  \TreePath \to \Maybe\naturals$ from \exampleref{tree example} to a
  dependent function $\types \lookup : \DFun{\BoundV t}{\Tree}
  \Fun{\NodeIn{\BoundV t}} \naturals$, as the path now must point to a node
  in the tree. We can define it as:
  \begin{gather*}
    \lookup \Big({\NodeV \ell n r}\Big) \Nil \definedby n
    \qquad
    \lookup \Big(\NodeV {t_{\Lft}}n {t_{\Rgt}}\Big) (d \:: p) \definedby \lookup \, t_{d} \, p
  \end{gather*}
  Note that since we have two trees $t_{\Lft}, t_{\Rgt}$ rather than a
  function $\TreeLabel \to \Tree$, the parametrisation $t_d$ that
  we use in the second case is just repetition in disguise. Later, we
  will show how to properly avoid this repetition by using the \WLOG
  construction.
\end{example}

\subsection{Indexed actions}
\subseclabel{indexed actions}

Just as in \exampleref{tree example}, the function $\lookup$ from
\exampleref{node lookup} is invariant with respect to flipping:
flipping both the tree and the path returns the same value. But note
that once we flip the tree $t$, the flipped path no longer lies in
$\NodeIn{t}$, but in $\NodeIn{(\swap \act t)}$. For example:
\[
\swap \act \Bigg(\Lft \:: \Rgt \:: \Nil \in \NodeIn   \bigg(
  \begin{tree diagram}
    \& \& 3 \\
    \& 2 \& \\
    1 \& \& 4
    \treeLink1322
    \treeLink2231
    \treeLink2233
\end{tree diagram}\bigg)\Bigg)
=
\Rgt \:: \Lft \:: \Nil \in \NodeIn   \bigg(
  \begin{tree diagram}
    3 \& \& \\
    \& 2 \& \\
    4 \& \& 1
    \treeLink1122
    \treeLink2231
    \treeLink2233
\end{tree diagram}\bigg)
\tag*{\qedhere}
\]
In general, an action over a family $\Ctx \types A$ should thus
consist of an action over the index set $\Ctx$, and an action over
each fibre $A \ctx$, which tracks the indexing action: if $a \in A
\ctx$, then $g \act a \in A (g \act \ctx)$.
%
\begin{definition}
  Given a $\G$-action $\Ctx$, an \emph{indexed $\G$-action $\Ctx
  \types \A$} consists of a \emph{carrier} family $\carrier\Ctx \types
  \carrier\A$ and a dependent function $(\iact{-}) : \DFun {\BoundV
  g}{\carrier \G} \DFun {\BoundV \ctx}{\carrier \Ctx}
  \Fun{\carrier\A\ctx} \carrier \A (g \act[\Ctx] \ctx)$, which we
  apply as $g \iact{\ctx} \ma$, such that
  \(
    \unit \iact{\ctx} \ma = \ma
  \) and \(
    (g \mul h) \iact{\ctx} \ma = g \iact{h \act \ctx} (h \iact{\ctx} \ma)
  \).
\end{definition}
%
When $\A$ is ambiguous, we add a superscript and write
$\iact[\A]{\ctx}$, and when  we can infer the index $\ctx$, we elide
it and write $g \act \ma$ instead of $g \iact{\ctx} \ma$.

\begin{example}
  The action of the cyclic group $\Cyc 4$ on 2048 boards can be
  generalised to an indexed action on the family of arbitrary matrices
  $\naturals \times \naturals \types \TConst{Mat}$. First, take $\Ctx$
  to be the action of $\Cyc 4$ on dimensions $\carrier\Ctx \definedby
  \naturals \times \naturals$, defined as:
  \(
    \Rotation\act\pair mn \definedby \pair nm
  \).
  We then have a $\Ctx$-indexed action defined as in \exampleref{2048
  action}:
  \[
    \Rotation
    \iact{\pair mn}
    \Bigg(
      \bigg(
    \begin{smallmatrix}
      x_{1,1} & \cdots & x_{1,n} \\
      \vdots & \ddots & \vdots \\[1ex]
      x_{m,1} & \cdots & x_{m,n}
    \end{smallmatrix}
    \bigg)
    \in
    \Mat mn\Bigg)
    =
    \begin{psmallmatrix}
      x_{1,n} & \cdots & x_{m,n} \\
      \vdots & \ddots & \vdots \\[1ex]
      x_{1,1} & \cdots & x_{m,1}
    \end{psmallmatrix}
    \in \Mat nm
    \tag*{\qedhere}
  \]
\end{example}

\begin{example}
  Take indexed $\G$-actions $\Ctx \types \A$ and $\Ctx, \x : \A \types
  \B$. Then, we can define the \emph{dependent pair} action $\Ctx
  \types \DPair\x\A\B$ as:
  \[
    \carrier{\Ctx \types \DPair\x\A\B} \definedby \carrier\Ctx \types \DPair\x{\carrier\A}{\carrier\B}
    \qquad
    \g \iact[\DPair\x\A\B]{\ctx}{\pair \ma\mb} \definedby
    \pair{\g\iact[\A]\ctx\ma}{\g\iact[\B]{\pair\ctx\ma}\mb}
  \]
  Furthermore, we can define the \emph{dependent function} action
  $\Ctx \types \DFun\x\A\B$ as:
  \begin{gather*}
    \carrier{\Ctx \types \DFun\x\A\B} \definedby \carrier\Ctx \types \DFun\x{\carrier\A}{\carrier\B}
    \\
    g\iact[\DFun\x\A\B]\ctx f \definedby
    \fun{(a : \carrier\A(g \act \ctx))}
      g \iact[\B]{\pair\ctx{\inv g \act\ma}} f_{\ctx}(\inv g \iact[\A]{g \act \ctx}\ma)
  \end{gather*}
  It is straightforward to check that both definitions satisfy the
  indexed action laws.
\end{example}

\begin{definition}
  Given a $\G$-action $\Ctx$ and an indexed $\G$-action $\Ctx \types
  \A$, a term $\carrier\Ctx \types t : \carrier\A$ is
  \emph{equivariant} when its components are: $t (g \act[\Ctx] \ctx) =
  g \iact{\ctx} t \ctx$ for all $\ctx \in \carrier\Ctx$, $g \in
  \carrier\G$.
\end{definition}

Equivariance of terms is preserved by reindexing and substitution.
Furthermore, an equivariant term $\Ctx, \x : \A \types t : \B$ induces
an equivariant dependent function $\Ctx \types (\fun{\x:\A}t) :
\DFun\x\A\B$, and two equivariant terms $\Ctx \types t : \A$ and $\ctx
: \Ctx \types s : \B \pair\ctx{t \ctx}$ induce an equivariant term
given by the dependent pair in each fibre
$\Ctx \types \pair ts : \DPair\x\A\B$. For more details and proofs,
see~\appref[C]{dependent types}.

The ease of compositionality should come as no surprise given that the
definition of an indexed action is in fact just a reformulation of the
notion of a family internally to the topos of $\G$-actions:

\begin{proposition}\propositionlabel{Grothendieck action}
  Let $\Ctx$ be a $\G$-action. Then, a $\Ctx$-indexed action $\A$
  amounts to an action $\coprod_{\Ctx}\A$ over the Grothendieck
  construction for which the projection is equivariant:
  \(
    \carrier{\coprod_{\Ctx}\A} = \coprod_{\carrier\Ctx}\carrier\A
  \) and \(
    \projection_1 :  \coprod_{\Ctx}\A \equito \Ctx
  \).
\end{proposition}

We have used this connection above to reformulate standard notions
from dependent types to the symmetric setting, and continue to do so
with canonisers, core functions, and the \WLOG construction. Ignoring
the additional indices, we recover the simply-typed notions from
\secref{wlog construction}.
%
\begin{definition}
  Let $\Ctx \vdash \A$ be an indexed $\G$-action. An \emph{indexed
  canoniser} for $\A$ is a $\carrier\Ctx$-indexed term $\carrier\Ctx \types \canoniser :
  \carrier\A \to \carrier\G$. For each $\ma \in \carrier\A\ctx$, we
  call the symmetry $\canoniser_{\ctx} \ma \in \carrier\G$
  the \emph{$\canoniser$-canonising symmetry} of $\ma$, and the
  element $\canoniser \ma\iact{\ctx} \ma$ the \emph{canonical form} of
  $\ma$. We define the family of \emph{canonical} elements
  $\carrier\Ctx \types \Canon\canoniser$ with the fibres
  \(
    (\Canon\canoniser) \ctx \definedby \set{
      \canoniser \ma\iact{\ctx'} \ma
      \suchthat
      \ctx' \in \Ctx,
      \ma \in \carrier\A\ctx',
      \canoniser \ma\act\ctx' = \ctx
    }
  \)
  for each
  $\ctx \in \carrier\Ctx$.
\end{definition}

\begin{example}
  To capture the symmetry of the lookup function from \exampleref{node
  lookup}, we define a canoniser for the indexed $\Sym 2$-action on
  $\mtree : \Tree \types \NodeIn\mtree$:
  \[
    \canoniser_{\Node \triple\ell n r}\Nil \definedby \stet \qquad
    \canoniser_{\Node \triple\ell n r}(\Lft \:: p) \definedby \stet \qquad
    \canoniser_{\Node \triple\ell n r}(\Rgt \:: p) \definedby \swap
  \]
  The canonical elements in $\NodeIn\btree$ are \( \Nil \) and \(
  (\Lft \:: p) \in \NodeIn(\Node\triple\ell n r) \) where $p \in
  \NodeIn\ell$ can be an arbitrary path, even one that moves right on
  any of its steps.
\end{example}

\begin{definition}
  Take indexed $\G$-actions $\Ctx \types \A$ and $\Ctx, x : \A \types
  \B$, a subfamily $\carrier\Ctx \types X \subseteq \carrier\A$, and
  a function $\carrier\Ctx \types f : \DFun{\x}{X}\carrier\B$.
  We write $\Ctx \types f : \EFun{\x}{X \subseteq \A}\B$ and say that
  $f$ is \emph{equivariant} when
  \(
    g \iact[\A]{\ctx} x \in X (g \act \ctx)
    \implies
    f_{g\act \ctx}(g\iact[\A]{\ctx}x) = g\iact[\B]{\pair\ctx x}f_{\ctx}x
  \)
  for all $g \in \carrier\G$, $\ctx \in \carrier\Ctx$, and $x \in X\ctx$.
\end{definition}
We can now define the dependently-typed version of the \WLOG{} construction:
\begin{definition}
  Let $\Ctx \types \A$ and $\Ctx, \x : \A \types \B$ be indexed
  $\G$-actions, $\Ctx \types \canoniser : \carrier\A \to \carrier\G$
  an indexed canoniser, and $\Ctx \types f :
  \DFun{\x}{\Canon\canoniser} \carrier\B$ a dependent core function.
  Then, the \emph{indexed \WLOG construction}~$\wlg(c,f) :
  \DFun{\x}{\carrier\A} \carrier\B$ is defined as:
  \[
  \myfigbox{$
    \wlg(\canoniser,f)_{\ctx} \definedby \fun{(\ma : \carrier\A\ctx)}
    \inv*{\canoniser_{\ctx}\ma}\iact[\B]{\pair{\canoniser_{\ctx} a \act[\Ctx] \ctx}{\canoniser_{\ctx} a \iact[\A]{\ctx} a}}
  f_{\canoniser_{\ctx}\ma \act[\Ctx] \ctx}(\canoniser_{\ctx} \ma\iact[\A]{\ctx} \ma)
  $}
  \]
\end{definition}

With all the indices precisely spelt out, the definition can look a
bit intimidating, but if we elide them we get
\(
  \wlg(\canoniser,f)
    = \fun{\ma} \inv*{\canoniser\ma}\iact[\B]{} f(\canoniser \ma\iact[\A]{} \ma)
\)
as in the simply-typed case. And just as before in
\theoremref{wlog-construction}, the \WLOG construction gives the
unique equivariant extension of a core function.

\begin{theorem}[indexed representation]
  \theoremlabel{indexed-wlog-construction}
  Let $\Ctx \types \A$ and $\Ctx, \x : \A \types \B$ be indexed
  $\G$-actions and $\Ctx \types \canoniser : \carrier\A \to
  \carrier\G$ an indexed canoniser for $\A$. If $\Ctx \types f :
  \EFun{\x}{\Canon\canoniser \subseteq \A} \B$ is an equivariant
  dependent core function, then  $\wlg(c,f) : \EFun{\x}{\A} \B$ is
  the unique equivariant function extending~$f$ to the whole
  $\carrier\Ctx \types \carrier\A$.
\end{theorem}

\begin{example}
  \examplelabel{pigeonhole principle}
  Let us revisit the pigeonhole principle~\cite{wiki:wlog} from the
  introduction and extend it in two ways. First, instead of a proof,
  let us construct a function
  \(
    \mathsf{findSame} : \DFun{\BoundV{xyz}}{\set{\RED, \BLU}^3} \Type{Same}{\BoundV{xyz}}
  \)
  where
  \(
    \Type{Same}{\BoundV{xyz}} \definedby \set{
      (i,j) : (\Fin 3)^2 \suchthat i \neq j \land \projection_i \BoundV{xyz} = \projection_j \BoundV{xyz}
    }
  \)
  is the family of distinct positions of the same colour, indexed by
  the colouring $\BoundV{xyz}$.

  And second, we shall see that in the special case of the proof,
  there is another opportunity to use the \WLOG{} argument.
  Informally, our proof goes as follows:
  %
  \begin{quote}
  \wLOG{}, the first object is red. The other two objects either have
  the same colour or different colours. If they have the same colour,
  choose $(1,2)$. If they have different colours, \WLOG{} the second
  object is red, so choose $(0,1)$.
  \end{quote}
  %
  Such nested uses of \WLOG{} may be suspicious and confusing, which
  is exactly why it is important to make all aspects of \WLOG{}
  arguments explicit. In this case, this use is valid. For example,
  when the second object is not red and the third is red, the
  inner argument exchanges their positions:
  \[
    \newcommand{\balls}[2]{\vcenter{\hbox{
      \begin{tikzpicture}[
        scale=0.5,
        obj/.style={circle, draw, inner sep=0pt, minimum size=4mm},
        rdot/.style={obj, fill=red!60, draw=red},
        bdot/.style={obj, fill=blue!20, draw=blue},
        pair/.style={rounded corners, thin, draw=gray},
        >=stealth, font=\small
      ]
        \pgfmathsetmacro\dx{-0.03}
        \foreach \col/\f in {#1}
          \node[\f] at (\col+0.5, 0.5) {\ifthenelse{\equal{\f}{bdot}}{b}{r}};
        #2
      \end{tikzpicture}
    }}}
    \balls{0/rdot, 1/bdot, 2/rdot}{}
    \qquad
    \explain{\rightleftharpoons}{switch\\[-5pt]objects}
    \qquad
    \balls{0/rdot, 1/rdot, 2/bdot}{}
    \qquad
    \explain{\mapsto}{pick the first\\[-5pt]and second object}
    \qquad
    \balls{0/rdot, 1/rdot, 2/bdot}{
      \draw[pair] (0-\dx,0-\dx) rectangle (1+\dx,1+\dx);
      \draw[pair] (1-\dx,0-\dx) rectangle (2+\dx,1+\dx);
    }
    \qquad
    \explain{\leftrightharpoons}{switch\\[-5pt]objects back}
    \qquad
    \balls{0/rdot, 1/bdot, 2/rdot}{
      \draw[pair] (0-\dx,0-\dx) rectangle (1+\dx,1+\dx);
      \draw[pair] (2-\dx,0-\dx) rectangle (3+\dx,1+\dx);
    }
  \]

  In the outer level, $\Sym 2$ acts on $\set{\RED,\BLU}$ by flipping
  the two colours, and on the domain $\set{\RED,\BLU}^3$ with the
  component-wise product action. On the codomain $\set{\RED,\BLU}^3
  \types \Type{Same}$, we lift the trivial action to preserve the
  elements of the fibre, but change the index. This is a sub-action
  (i.e. an action that preserves a subset) since if objects at two
  positions are of the same colour, they remain so after we exchange
  the colours. For example, $\smash{\swap\iact{(\RED,\RED,\BLU)}(0,1)}
  = (0,1) \in \Type{Same}{(\BLU, \BLU,\RED)}$

  By swapping colours, we make the first object red, thus the outer
  canoniser is:
  \[
    \mathsf{makeFirstRed}(\RED,\mb,\mc) \definedby \stet
    \qquad
    \mathsf{makeFirstRed}(\BLU,\mb,\mc) \definedby \swap
  \]
  and using it, we can define
  \[
  \mathsf{findSame} \definedby \wlg\Bigg(
    \mathsf{makeFirstRed},
    \fun{(\RED, \mb, \mc)}
      \begin{cases}
        (1,2) &(\mb = \mc) \\
        \mathsf{findOtherRed}(\mb,\mc) &(\mb \neq \mc)
      \end{cases}
  \Bigg)
  \]
  where we still need to define
  \[
  \mathsf{findOtherRed} : \DFun{(\y,\z)}{\set{(\mb,\mc) \in \set{\RED,\BLU}^2\suchthat \mb \ne \mc}}
  \Type{Same}{(\RED,\y,\z)}
  \]
  Again, we employ symmetry, but now, the $\Sym2$-action on the domain
  swaps the two components of the pair, but not their colour, i.e.,
  $\swap\act(\mb,\mc) \definedby (\mc, \mb)$. For the codomain, define
  an action on $\Fin 3$ that swaps $1$ and $2$, and extend it
  component-wise to the product $(\Fin 3)^2$. Since $\mb$ and $\mc$
  are of different colours, it is easy to check that this action again
  preserves the defining property of $\Type{Same}$ and thus lifts to a
  sub-action on the codomain. For example,
  $\smash{\swap\iact{(\RED,\BLU)}(0,1) = (0,2) \in \Type{Same}{(\RED, \BLU,
  \RED)}}$.

  The inner canoniser now makes the second object red:
  \[
  \mathsf{makeSecondRed}(\RED, \BLU) \definedby \stet
  \qquad
  \mathsf{makeSecondRed}(\BLU, \RED) \definedby \swap
  \]
  and so
  \(
  \mathsf{findOtherRed}\definedby \wlg(\mathsf{makeSecondRed},
  \fun{(\RED,\BLU)}(0,1))
  \).
  Equivariance is immediate since both core functions are vacuously
  equivariant.
\end{example}

\subsection{Propositional reasoning}
\subseclabel{propositional reasoning}

So far we have used the \WLOG{} construction to define functions. To
formalise mathematical proofs, we apply the Curry-Howard
correspondence and identify propositions with types, and their proofs
with terms. We could have also approached this differently and e.g.,
transferred the \WLOG{} construction to higher-order logic or
first-order logic, but this would require us to duplicate the
definitions for concepts such as group actions or equivariance for
proofs.

For concreteness, let us first borrow a running example of Schur's
inequality from \citet{DBLP:conf/tphol/Harrison09}, which states that
for all $x, y, z, t \ge 0$, we have
\[
  x^{t}(x - y)(x - z)+y^{t}(y - z)(y - x)+z^{t}(z - x)(z - y) \ge 0
\]
Since the inequality is symmetric in the three variables, the standard
proof assumes that \WLOG, $x \le y \le z$ and rearranges the
expression to an evidently non-negative form.

Let $\Prop \definedby \set{\Tru,\Fls}$ be the set of logical values.
Each \emph{proposition} $\Ctx \types \phi : \Prop$ determines a family
$\Ctx \types \Type{Proof} \phi$ such that for each $\ctx \in \Ctx$,
the fibre $(\Type{Proof} \phi) \ctx = \terminal = \set{\unitval}$ when
$\phi\ctx = \Tru$ and $(\Type{Proof} \phi) \ctx = \emptyset$ when
$\phi\ctx = \Fls$. As is conventional, we overload the notation and
identify propositions with their proofs, writing $\Ctx \types
\Type{Proof} \phi$ simply as $\Ctx \types \phi$. Then, any proof term
$\Ctx \types p : \phi$ is an explicit witness that $\phi$ holds at
every $\ctx \in \Ctx$.

\begin{example}
  \examplelabel{schur invariance}
  Schur's inequality is the proposition $(\x, \y, \z) :
  \preals^3 \types \Type{Schur}(x, y, z) : \Prop$ where
  \[
    \Type{Schur}(x, y,z) \definedby \forall t \ge 0.\
    x^{t}(x-y)(x-z)+y^{t}(y-z)(y-x)+z^{t}(z-x)(z-y) \ge 0
  \]
  Recall from \exampleref{symmetric-group-action} that $\Sym 3$ acts
  on $\preals^3$ by permutation. The proposition $\Type{Schur}$ is
  invariant under those permutations.
\end{example}

\begin{definition}
  A proposition $\Ctx \types \phi : \Prop$ is \emph{invariant} when it
  is equivariant as a term with respect to the trivial $\G$-action on
  $\Prop$, i.e. $\phi (\g \act \ctx) = \phi \ctx$ for all $\ctx \in
  \Ctx$.
\end{definition}

In addition to equivariant substitutions, like all equivariant terms,
invariant propositions are also preserved by logical connectives and
quantifiers. Each invariant proposition $\Ctx \types \phi : \Prop$
induces a $\Ctx$-indexed action on its proof family $\Ctx \types
\phi$. Moreover, given a $\Ctx$-indexed action $\A$, each invariant
proposition $\Ctx, \x : \A \types \phi : \Prop$ induces a
$\Ctx$-indexed action on the subfamily $\Ctx \types \set{\x : \A
\suchthat \phi} \subseteq \A$. For more details and proofs, again
see~\appref[C]{dependent types}.

\begin{example}
  Coming back to Schur's inequality, we can see that the \WLOG
  argument again fits the now familiar pattern:
  \[
    x_0, x_1, x_2 \in \preals
    \quad
    \explain\mapsto{sort using a\\[-5pt]permutation $\perm^{-1}$}
    \quad
    x_{\perm0} \le x_{\perm1} \le x_{\perm2}
    \quad
    \explain\mapsto{prove using \\[-5pt]additional assumption}
    \quad
    \Type{Schur}(x_{\perm0}, x_{\perm1}, x_{\perm2})
    \quad
    \explain\mapsto{$\Type{Schur}$ is invariant\\[-5pt]under permutations}
    \quad
    \Type{Schur}(x_0, x_1, x_2)
  \]
  Let us use the \WLOG construction to obtain a proof, i.e. a
  function $\DFun{(\x, \y, \z)}{\preals^3} \Type{Schur}(x, y, z)$.
  From \exampleref{schur invariance} we already have an $\Sym
  3$-action on $\preals^3$, and since $\Type{Schur}$ is an invariant
  proposition, we also get an indexed $\Sym 3$-action on $\preals^3
  \types \Type{Schur}$. For the canoniser, we take a stable sorting
  function $\sort : \preals^3 \to \Sym 3$ from \exampleref{sorting
  canoniser} so that $\Canon\sort$ is exactly the set of sorted
  triples. The last thing we need is the core function $f :
  \set{(x, y, z) : \preals^3 \suchthat x \le y \le z} \to \Schur(x, y,
  z)$, which we can define using the rearrangement in the standard proof.
\end{example}

The resulting function $\wlg(\sort, f)$ is equivariant since in the
case of propositions, core functions are always equivariant due to
proof irrelevance:

\begin{lemma}\lemmalabel{core functions to propositions invariance}
  Let $\Ctx, \x : \A \types \phi : \Prop$ be an invariant proposition,
  and $\Ctx \types \canoniser : \carrier\A \to \carrier\G$ a
  canoniser.  Every core function $\Ctx \types f : \Canon\canoniser
  \to \phi$ is equivariant.
\end{lemma}

Harrison suggests proving Schur's inequality using the following
tactic\footnotemark{}, which states that it
is sufficient to prove permutation invariant propositions only for
sorted triples~\cite[page 44]{DBLP:conf/tphol/Harrison09}:
\begin{align*}
  \inferrule{
    \forall \x,\y,\z .\ P(\x,\y,\z) \Rightarrow P(\y,\x,\z) \land P(\x,\z,\y)\\
    \forall \x,\y,\z .\ (\x \le \y) \land (\y \le \z) \Rightarrow P(\x,\y,\z)
  }{
    \forall \x,\y,\z .\ P(\x,\y,\z)
  }
\end{align*}
\footnotetext{Available here: \url{https://github.com/jrh13/hol-light/blob/e5236a5/int.ml\#L612}.}%
%
Let us see how this bespoke tactic relates to our general \WLOG{}
construction. In contrast to our approach, the symmetry group $\Sym 3$
is left implicit, as is the canoniser. Instead, only its set of
canonical elements $\set{(x, y, z) : \preals^3 \suchthat x \le y \le
z}$ is spelt out. Furthermore, the first assumption corresponds to an
indexed $\Sym 3$-action on the proposition $P$, and the second
assumption is exactly the core function.

The latter is evident, while for the former, first recall from
\exampleref{2048 action} that a group action is already uniquely
determined by the action of any set of generators. The symmetry
group $\Sym 3$ is generated by the two transpositions $(0
\leftrightarrow 1)$ and $(1 \leftrightarrow 2)$, hence the action of
$\Sym 3$ is determined by maps $(0 \leftrightarrow 1) \act {-}$ and
$(1 \leftrightarrow 2) \act {-}$. Spelling out the indices more
precisely, these two maps are:
\begin{align*}
  (0 \leftrightarrow 1) \iact{(x, y, z)} {-} &: P(x, y, z) \to P((0 \leftrightarrow 1) \act (x, y, z))
  &\text{equivalently}\quad
  &P(x, y, z) \to P(y, x, z)
  \\
  (1 \leftrightarrow 2) \iact{(x, y, z)} {-} &: P(x, y, z) \to P((1 \leftrightarrow 2) \act (x, y, z))
  &\text{equivalently}\quad
  &P(x, y, z) \to P(x, z, y)
\end{align*}
And this is exactly the contents of the first assumption.

As in \exampleref{2048 action}, we must ensure that such a definition
by generator cases respects all the equalities of the group. For
example in any $\Sym n$-action, $(0 \leftrightarrow 1) \act {-}$ must
be its own inverse (amongst other relations) since $(0 \leftrightarrow
1) \mul (0 \leftrightarrow 1) = \unit$. However, in the case of
propositions, once again proof irrelevance implies this condition.

Exposing the symmetric structure in proofs also allows us to use
existing results in group theory. Consider Harrison's proof in the
tactic, which explicitly lists all six
permutations~\cite{DBLP:conf/tphol/Harrison09}. For $\Sym 3$, this
number may still be manageable, but in general, the group~$\Sym n$ has
$n!$ elements, so when $n$ increases, the size of the proof script
quickly becomes prohibitively large. Also, such an approach precludes
a proof of a general $\wlg_n$ principle if $n$ is a variable, as one
cannot explicitly list all the $n!$ cases in the script. In contrast,
we can lean on the general group-theoretic result, whose proof is
independent of the concrete value of $n$, that $\Sym n$ is generated,
for example, by a single cycle $(0 \, 1 \, 2 \, \cdots \, n - 1)$ and
a single transposition of adjacent elements such as $(0
\leftrightarrow 1)$. So it suffices to prove two cases
to define any propositional $\Sym n$-action, yielding a single
parametric principle $\wlg_n$.

\subsection{Equivariant canonisers}
\subseclabel{equivariant canonisers}

In \subsecref{wlog-equivariance}, we saw that all \WLOG{} constructions
of a canoniser are equivariant iff the canoniser is strictly
equivariant (\theoremref{all core functions are partially
  equivariant}). This result is significant because it allows us to do
away completely with the equivariance of the core function, which
involves comparing every pair of symmetric canonical elements.
However, only the semi-regular actions can
have such canonisers~(\lemmaref{semi-regularity}). We now investigate a weaker
notion of equivariance. It is also characterised through equivariant
extension, but of \emph{stable} core functions. A core function $f$ is
stable when every \emph{stabilising} symmetry~$g$ of an
element---$g\act a = a$---also stabilises the image $f\,a$. This
necessary condition is natural and often easy to check.  Furthermore,
if we assume the axiom of choice, all actions admit equivariant
canonisers, making them a viable---albeit theoretical---class of canonisers.

Both $\canoniser (g \act
a)$ and $\canoniser a \mul \inv g$ canonise $g \act a$. Strictly
equivariant canonisers equate these two symmetries:
$\canoniser (g \act a) = \canoniser a \mul \inv g$.
Equivariant canonisers merely relate
them: $\canoniser (g \act a) \stabeq{} \canoniser a \mul \inv g$.

\begin{definition}
  \definitionlabel{indistinguishability}
  Take a $\G$-action $\A$. Then, for each $a \in \carrier\A$, we equip
  $\carrier\G$ with an \emph{indistinguishability} relation $g \stabeq
  a h \iff g \act a = h \act a$. Let $[g]_{\stabeq a}$ be the
  equivalence class of $g$ in $\carrier\G / \stabeq a$. Define an
  indexed $\G$-action $\A \types \G / \A$ with the fibre $(\G / \A) \, a
  \definedby \carrier\G / \stabeq a$ and with $\act[\G / \A]$ defined
  as:
  $
  g \iact[\G / \A] a{[h]_{\stabeq a}} \definedby [h \mul \inv g]_{\stabeq {g \act a}}
  $.
\end{definition}

The action is well-defined since $h_1 \stabeq a h_2$ implies $h_1 \mul
\inv g \stabeq{g \act a} h_2 \mul \inv g$. The relation $\stabeq a$ is
closely related to the \emph{stabiliser} $\stabilise a \definedby
\set{g \in \carrier\G \suchthat g \act a = a} = [\unit]_{\stabeq a}$, i.e. a subgroup of all
symmetries that preserve $a$. In particular, $g \stabeq a h$ is
equivalent to $(\inv h \mul g) \act a = a$, i.e. $\inv h \mul g \in
\stabilise a$.

\begin{remark}
The equivalence classes of the indistinguishability relation are
precisely the cosets of the stabiliser group. Therefore, we can recast the
celebrated Orbit-Stabiliser theorem as an \emph{equivariant} isomorphism between
the actions on orbits and stabiliser cosets.

Given a $\G$-action $\A$, the fibres of the dependent family of
\emph{orbits} contain all the symmetric elements of an index, i.e. $a
: \carrier\A \types \OrbitOf[\G]a \definedby \set{h \act a \suchthat h
\in \carrier\G}$. On orbits, we can define an indexed action as:
\[
  g \iact a (h \act a \in \OrbitOf a)
  \definedby
  h \act a
  =
  (h \mul \inv g) \act (g \act a)
  \in
  \OrbitOf{g \act a}
\]
Even though the action does not change the fibres, it is not
trivial, since it changes the indices.
\end{remark}

\begin{theorem}[Orbit-Stabiliser]
  \theoremlabel{orbit stabiliser}
  Let $\A$ be a $\G$-action. We have an equivariant isomorphism:
  \[
  a : \A \types \phi_a : (\G / \A) a \overset{\isomorphic}\equito \OrbitOf{a}
  \qquad
  \phi_a [h]_{\stabeq a} \definedby h \act a
  \qquad
  \inv \phi_a(h \act a) \definedby [h]_{\stabeq a}
  \]
\end{theorem}

The indexed action $\G / \A$ generalises $\inv\G$ and allows us to
consider a weaker property in which the equality $\canoniser\parent{g
\act[\A] a} = g \act[\inv\G] \canoniser a$ of strictly equivariant
canonisers is replaced by indistinguishability:

\begin{definition}
  A canoniser $\canoniser : \carrier\A \to \carrier\G$ is
  \emph{equivariant} when post-composing it with the quotient map
  yields an equivariant map $(\fun a[\canoniser\,a]_a) : \EFun
  a\A{(\carrier\G / \A) a}$, i.e. when $\canoniser (g \act[\A]
  a) \stabeq{g \act a} g \act[\inv\G] \canoniser a$.
\end{definition}

\begin{example}
  % \examplelabel{equivariant canonisers}
  For a totally ordered set $P$, we can define the $\Sym 2$-action on
  $P \times P$ by $\swap\act\pair pq \definedby \pair qp$ and a
  canoniser $\canoniser : P \times P \to \carrier{\Sym2}$ such that
  $\canoniser \pair pq = \stet$ if $p \le q$, and $\canoniser \pair pq
  = \swap$ if $p > q$. Then $\canoniser$ is not strictly equivariant
  since for any $p \in P$, we have:
  \[
  \canoniser\parent{\swap \act \pair pp} = \canoniser\pair pp = \stet \neq \swap
  = \stet \mul {\inv\swap} = \swap\iact{\pair pp}\canoniser\pair pp
  \]
  It is, however, equivariant, since
  \(
  \canoniser\parent{\swap \act \pair pp} =
  \stet \stabeq{\pair pp} \swap =
  \swap\iact{\pair pp}\canoniser\pair pp
  \).
\end{example}

Note that for every equivariant function $f$ and every stabilising
element we have:
\[
\begin{tabular}[c]{@{}c@{}}
\myfigbox{
  $
  g \in \stabilise a
  \implies
  g \act f a \explain={equivariance} f(g \act a) \explain{=}{stabiliser} f a
  \implies
  g \in \stabilise {f \, a}
  $
}
\end{tabular}
\qquad\text{i.e.:}\qquad
\myfigbox{$
  \stabilise a \subseteq \stabilise {f \, a}
$}
\]
We thus arrive at the following condition on core functions:
\begin{definition}
  Take indexed $\G$-actions $\Ctx \types \A$ and $\Ctx, x : \A \types
  \B$, a subfamily $\carrier\Ctx \types X \subseteq \carrier\A$. A
  function $\carrier\Ctx \types f : \DFun{\x}{X}\carrier\B$
  is \emph{stable} when it preserves the stabiliser: $\stabilise x
  \subset \stabilise {f\,x}$ for all $x \in X\ctx$.
\end{definition}

\begin{example}
  Just as in \exampleref{strictly equivariant canonisers}, consider a
  non-total ordering $P$ and a stable sorting function, and the
  permutation sub-action on all tuples with distinct elements
  $\set{\vec p \in P^n \suchthat \forall i,j \in \Fin n. p_i \neq
  p_j}$. Then, we can define $\minimal : \carrier P^n \to \Fin n$ that
  finds the index of a minimal element by $\minimal \definedby \wlg(\sort,
  (\fun{\_}0))$.

  Now, for any two incomparable $p, q$, both $\pair pq$ and $\pair qp$
  are sorted, thus canonical. However:
  \[
    (0 \leftrightarrow 1) \act ((\fun{\_}0) \seq{p,q})
    = (0 \leftrightarrow 1) \act 0
    = 1 \neq 0 = (\fun{\_}0) \seq{q,p}
    = (\fun{\_}0) \parent{(0 \leftrightarrow 1) \act \seq{p,q}}
  \]
  and so $(\fun{\_}0)$ is not equivariant on $\Canon\sort$. It is,
  however, vacuously stable. Indeed, take a canonical $\vec p$ and a
  $\vec p$-stabilising permutation $\perm \in \stabilise{\vec p}$.
  Since all the elements in $\vec p$ are distinct, yet $\perm\act \vec
  p = \vec p$, we have that $\perm = \id$. Therefore $\perm \in
  \stabilise{f\vec p} = \stabilise 0$.
\end{example}

So, in general, stability is strictly weaker than equivariance.
Nonetheless, in the presence of an equivariant canoniser, stability
and equivariance coincide.

\begin{theorem}\theoremlabel{all stable core functions are equivariant}
  Let $\Ctx \types \A$ be an indexed $G$-action and let $\carrier\Ctx \types
  \canoniser : \carrier\A \to \carrier\G$ be an indexed canoniser.
  Then, $\canoniser$ is equivariant iff for every indexed $G$-action
  $\Ctx, x : \A \types \B$, every stable $\Ctx$-indexed core function
  $\Ctx \types f : (x : \Canon\canoniser) \to \carrier\B$
  is equivariant.
\end{theorem}

Unlike before, where only semi-regular actions admitted strictly
equivariant canonisers, here every action admits an equivariant
canoniser, although we may need to invoke the axiom of choice.

\begin{theorem}\theoremlabel{equivariant canoniser choice}
  Using the axiom of choice, every indexed $\G$-action has an equivariant canoniser.
\end{theorem}

To us, this use of the axiom of choice indicates that identifying
symmetry is a creative activity, rather than a mechanical one.
